\chapter*{Abstract}
\addcontentsline{toc}{chapter}{Abstract}

In dieser Arbeit wird ein fiktives Auftragsszenario
bearbeitet, welches an den Versicherungskonzern Allianz SE angelehnt ist.
Im Speziellen beschäftigen die Allianz in diesem Szenario
zwei Anwendungsfälle, die datengetrieben analysiert und gelöst
werden sollen.
Erstens die Erkennung von Betrug bei Kfz-Schadensmeldungen.
Und zweitens die Prognose kündigungsgefährdeter 
Krankenversicherungs-Kunden und -Kundinnen.
Es wird angenommen, dass die Allianz noch keine
Infrastruktur für die analytische Verarbeitung von Daten betreibt.
In dieser Arbeit wird eine umfassende Datenarchitektur sowie 
eine beispielhafte Modellierung und Aufbereitungspipeline vorgestellt,
mit deren Hilfe die beiden Anwendungsfälle abgedeckt werden.
Sowohl die Auswahl der Datenarchitekturkomponenten als auch 
die Abfolge der Datenaufbereitungsschritte sind auf 
die Rahmenbedingungen der beiden Use Cases und der Allianz als Ganzes abgestimmt.
Basierend auf dieser Ausarbeitung könnte in der Folge
eine passende technische Infrastruktur aufgebaut werden.
